\section{System Architecture}


Hanzo Network implements a layered architecture separating consensus, compute orchestration, and AI inference concerns while maintaining tight integration for performance.

\subsection{Core Components}

\subsubsection{Node Architecture}

The Hanzo node consists of 25+ Rust crates organized into functional layers:

\begin{itemize}
\item \textbf{Consensus Layer} (\texttt{hanzo-consensus}): Lux Wave consensus with sled database backend
\item \textbf{Runtime Layer} (\texttt{hanzo-runtime}): Unified abstraction for Docker/K8s/WASM/Firecracker
\item \textbf{Cryptography Layer} (\texttt{hanzo-pqc}, \texttt{hanzo-kbs}): Post-quantum KEM/DSA with Key Broker Service
\item \textbf{AI Layer} (\texttt{hanzo-hllm}, \texttt{hanzo-embedding}): Routing and inference coordination
\item \textbf{Network Layer} (\texttt{hanzo-libp2p-relayer}): P2P networking with relay support
\item \textbf{API Layer} (\texttt{hanzo-http-api}): REST/WebSocket/SSE endpoints
\end{itemize}

\subsubsection{Port Configuration}

\begin{itemize}
\item \textbf{3690}: Main HTTP/WebSocket API
\item \textbf{3691}: P2P networking port
\item \textbf{3692}: Server-Sent Events (SSE) streaming
\item \textbf{36900}: Hanzo Engine inference service
\item \textbf{9650-9651}: Consensus RPC endpoints
\end{itemize}

